\documentclass[12pt, a4paper]{article}
\usepackage[utf8]{inputenc}
\usepackage[T1]{fontenc}
\usepackage{lmodern}
\usepackage{amsmath}
\DeclareUnicodeCharacter{2212}{\ensuremath{-}}
\DeclareUnicodeCharacter{00D7}{\ensuremath{\times}}
\DeclareUnicodeCharacter{2192}{\ensuremath{\rightarrow}}
\DeclareUnicodeCharacter{0394}{\ensuremath{\Delta}}
\DeclareUnicodeCharacter{03B2}{\ensuremath{\beta}}
\DeclareUnicodeCharacter{03B7}{\ensuremath{\eta}}
\DeclareUnicodeCharacter{03B1}{\ensuremath{\alpha}}
\DeclareUnicodeCharacter{03B8}{\ensuremath{\theta}}
\DeclareUnicodeCharacter{03C4}{\ensuremath{\tau}}
\DeclareUnicodeCharacter{03C3}{\ensuremath{\sigma}}
\DeclareUnicodeCharacter{2264}{\ensuremath{\leq}}
\DeclareUnicodeCharacter{2208}{\ensuremath{\in}}
\DeclareUnicodeCharacter{2260}{\ensuremath{\neq}}
\DeclareUnicodeCharacter{00B1}{\ensuremath{\pm}}
\DeclareUnicodeCharacter{00B2}{\textsuperscript{2}}
\DeclareUnicodeCharacter{2081}{\textsubscript{1}}
\DeclareUnicodeCharacter{2082}{\textsubscript{2}}
\DeclareUnicodeCharacter{25B3}{\ensuremath{\triangle}}
\DeclareUnicodeCharacter{2032}{\ensuremath{\prime}}
\DeclareUnicodeCharacter{0304}{\={}}
\usepackage{booktabs}
\usepackage{longtable}
\usepackage{array}
\usepackage{calc}
\usepackage{float}
\usepackage{graphicx}
\PassOptionsToPackage{hyphens}{url}
\usepackage{hyperref}
\usepackage{xurl}
\usepackage{setspace}
\usepackage[margin=1in]{geometry}
\usepackage[font=it,justification=centering]{caption}
\usepackage{times}

\setstretch{2.0}
\renewcommand{\arraystretch}{1.6}
\urlstyle{same}
\setlength{\emergencystretch}{3em}
\newenvironment{referenceshang}{%
  \par\begingroup
  \setlength{\parindent}{0pt}%
  \setlength{\parskip}{0.5\baselineskip}%
  \everypar{\hangindent=1.5em\hangafter=1\relax}%
  \ignorespaces
}{%
  \par\endgroup
}

\renewcommand{\thetable}{\Roman{table}}
\begin{document}

% ============================================================================
% JPBM Round 2 working draft, v3 (2026-08-26).
% v3 policy: JPBM is a fresh submission; the main text carries only required
% disclosures (preregistration deviation, measured mediators, factual condition
% description, non-equivalence, observational field data, Limitations).
% All other qualifiers live in Supplementary Information H and the Limitations
% section. Strengths are stated in the main text.
% All statistics verified against the released data files (04_analysis/verification).
% Anonymized: no author-identifying information.
% ============================================================================

{\centering\Large\bfseries
How Launch Terms Shape Failure Judgments\par
from Public Market Outcomes\par
in Digital Luxury Brand Extensions\par}

\section*{Structured Abstract}

\textbf{Purpose:} This research examines how a luxury brand's launch terms for a digital extension --- its initial price or free distribution --- shape the extension's initial evaluation, the Failure Inference formed once displayed buyer and market-value information appears, and Loss of Brand Luxuriousness.

\textbf{Design/methodology/approach:} A descriptive field study of 65 brand NFT collections and 22,841 X posts provides observational context. Six experiments (\emph{N} = 1,917), five preregistered and one supplementary, test the predictions; Studies 3a--3c hold displayed outcome information identical across conditions.

\textbf{Findings:} Before buyer information appeared, a low-priced extension produced higher Loss of Brand Luxuriousness than a free one. A complimentary condition raised Loss of Product Luxuriousness for a physical but not a digital offering. Under identical zero-buyer information, complimentary or free conditions lowered Failure Inference, with a negative indirect effect on Loss of Brand Luxuriousness; the reduction weakened for a central-logo image and an official-website source, and one study showed inconsistent mediation.

\textbf{Research limitations/implications:} Indirect effects are process-consistent estimates tied to the implemented conditions; the field evidence is observational.

\textbf{Practical implications:} A low launch price is not necessarily safer for the parent brand than free distribution; launch terms should be planned jointly with the market information platforms later display.

\textbf{Originality:} Attribution theory connects consumers' interpretation of launch terms with their interpretation of the same extension's subsequently displayed market outcomes, separating these judgments from parent-brand feedback effects.

\emph{Keywords}: digital brand extension, brand extension feedback effects, attribution theory, zero-price effect, luxury brand, NFT

\emph{Paper type}: Research paper

\newpage

\section{Introduction}

% I1: The phenomenon of digital brand extensions
Luxury brands increasingly introduce digital brand extensions in the form of branded collectibles that may be sold or distributed free of charge and later traded (Bao \emph{et al.}, 2025; Brandes and D\"olp, 2025). Recent launches such as Prada's Timecapsule and Louis Vuitton's digital collectibles illustrate the format (Prada, 2026; Louis Vuitton, 2025): the offering carries the brand's name, logo, and design; it reaches consumers through an initial sale or free distribution; and it can subsequently be traded on marketplaces. In this research, non-fungible tokens (NFTs) serve as the empirical setting for such extensions rather than as a theoretical construct in their own right. Because these offerings remain visibly connected to the parent brand, information about the extension can also affect evaluations of the parent brand.

% I2: Spillover to parent-brand evaluation
A distinctive feature of this setting is that marketplaces display the extension's market outcomes. Displayed information showing no buyers, inactive trading, or little market value can serve as an input to consumers' evaluations of the parent brand: outcomes such as Porsche's halted 911 NFT sale and the litigation that followed Nike's closure of its NFT platform illustrate how visibly weak outcomes can accompany harm to the brand (CoinDesk, 2023; The Star, 2025). Research on brand extension feedback effects shows that extension information can alter parent-brand evaluations (Ahluwalia and G\"urhan-Canli, 2000; Milberg \emph{et al.}, 2023), and the accessibility--diagnosticity perspective explains why accessible extension information may be used in those judgments (Ahluwalia and G\"urhan-Canli, 2000). Our focal outcome is the three-item measure \emph{Loss of Brand Luxuriousness}, assessed immediately after exposure; the research therefore concerns short-run perceptions rather than long-term brand equity. Consumers also see how the extension was initially priced or distributed, and that information may shape how they interpret the later buyer and price information.

% I3: What existing research explains
Existing research explains separate parts of this sequence but has not jointly examined how initial price or distribution information affects responses to subsequently displayed buyer and market-value information. Attribution theory holds that consumers infer causes, motives, and meanings from observable firm actions (Folkes, 1988; Weiner, 2000), including pricing actions (Campbell, 1999). Pricing research shows that a low positive price and a zero price do not carry the same meaning (Shampanier \emph{et al.}, 2007; Palmeira and Srivastava, 2013). Research on responses to unfavorable outcomes shows that judgments about cause, controllability, and responsibility change consumer reactions (Folkes, 1984; Weiner, 2000). Finally, research on brand extension feedback effects shows that extension information can spill over to parent-brand evaluation (G\"urhan-Canli and Maheswaran, 1998; Ahluwalia and G\"urhan-Canli, 2000). What remains unclear is whether initial price or distribution information affects both the initial evaluation of a digital brand extension and the later Failure Inference formed after consumers see buyer and market-value information.

% I4: Two distinct judgments
Initial price or distribution information may accordingly affect two distinct consumer judgments. The first is evaluative: how should the offering itself be evaluated? The second arises only after outcome information appears: after seeing the displayed buyer count and market value, does the launch appear regrettable, poorly judged, or problematic for the company? The same price or distribution information may shape both the initial evaluation of the extension and the later judgment captured by Failure Inference. We examine these judgments separately and then assess their immediate consequences for perceived brand luxuriousness. This two-judgment sequence mirrors the architecture of crisis research, in which firm actions taken before any adversity shape how consumers respond to a later negative event (Dawar and Pillutla, 2000; Klein and Dawar, 2004); here the prior action is the launch itself.

% I5: Initial evaluation of the offering
Before any buyer-count or market-value information is shown, price or distribution information can affect how consumers value the extension. Study 1 compares four price conditions for a Gucci NFT and tests whether a low positive price, which presents the extension as a low-valued paid offer, produces higher Loss of Brand Luxuriousness than free distribution; price--quality inference motivates this contrast (Rao and Monroe, 1989). Study 2 then contrasts a comparable-price condition with a complimentary-with-a-fashion-magazine condition for physical and digital versions of the same offering. The magazine-gift condition is a distinct implemented condition from Study 1's plain free distribution, and we analyze it as such (Table II). Study 3a then asks whether Failure Inference differs between the purchased/comparable and complimentary magazine-gift conditions after identical buyer and market-value information is shown.

% I6: Displayed buyer information and Failure Inference
Information showing ``Buyers: 0'' and almost no market value is unfavorable, but it does not by itself establish that consumers will judge the initiative as a failure. The badness of a displayed outcome and the judgment that the launch was a failure are distinct: we capture the latter with \emph{Failure Inference}, a three-item measure assessing whether the launch appears regrettable, a bad decision, or a serious problem for the company. The purchased/comparable condition presents the initial distribution as a paid market exchange, whereas the complimentary magazine-gift condition does not present it as the same paid exchange; the two conditions may therefore give the same displayed outcome different meanings. Study 3a therefore tests whether Failure Inference differs between the two conditions under the same displayed buyer and market-value information.

% I7: Studies 3b and 3c as concrete presentation contrasts
Studies 3b and 3c test whether the Study 3a pattern changes across two concrete presentation contrasts, rather than treating them as indicators of a single moderator. Study 3b varies the extension's image: an image dominated by a central GG logo versus a bee image with a small logo, a contrast informed by brand prominence research (Han \emph{et al.}, 2010). Study 3c varies where the launch and its outcome information are presented: an official brand website versus a news article. We therefore report separately the contrast between the image with a central GG logo and the bee image with a small logo, and the contrast between the official brand website and the news article.

% I8: Research program and contributions
A Preliminary Field Study first describes associations between a Free/Paid classification of brand NFT collections and observed marketplace indicators, and between that classification and the VADER compound scores of X posts retrieved with collection-specific queries. Study 1 tests H1 (low versus free); Study 2 tests H2 (physical/digital $\times$ comparable price/complimentary magazine gift); Studies 3a--3c test H3a--H3c under displayed information showing zero buyers and almost no market value; a further preregistered study on favorable displayed information (Study 4) is reported as supplementary material. Together, the program combines two field datasets with six experiments (\emph{N} = 1,917), five of them preregistered. The research makes three contributions. First, it separates outcome valence from failure judgment: displayed information showing no buyers is unfavourable in every condition, yet whether consumers read it as a failed launch varies with how the extension was launched. Second, it identifies launch terms as interpretive priors, showing that price and distribution information chosen before any outcome is observable alters how diagnostic objectively identical market information later appears. Third, it adds a pre-outcome determinant to research on brand extension feedback effects, in which how far an extension's market outcome travels to the parent brand depends on the terms under which it was launched as well as on the outcome itself. Throughout, the theoretical attribution processes are distinguished from the judgments directly measured, and the free and complimentary conditions are analyzed as the concrete implementations actually used, separating their effects on initial valuation, Failure Inference, and immediate Loss of Brand Luxuriousness. Table I sets out the key terms and the scope in which each is used here.

\begin{table}[H]
\centering
\caption{Key Terms and Their Scope in This Research}
\begingroup
\doublespacing
\renewcommand{\arraystretch}{1.0}
\begin{tabular}{@{}
  >{\raggedright\arraybackslash}p{(\columnwidth - 2\tabcolsep) * \real{0.30}}
  >{\raggedright\arraybackslash}p{(\columnwidth - 2\tabcolsep) * \real{0.70}}@{}}
\toprule
\textbf{Term} & \textbf{Definition (scope)} \\
\midrule
Digital brand extension & A brand-issued digital collectible carrying the parent brand's name, logo, or design; implemented here as NFTs that are transferable on secondary markets (Sung \emph{et al.}, 2023; Brandes and D\"olp, 2025). \\
Initial price and distribution information & The price or distribution terms shown at launch (e.g., high, comparable, or low price; free distribution; a complimentary gift with a magazine or fan book). Exact wording per study appears in Table II. \\
Displayed buyer and market-value information & Stimulus content showing the extension's marketplace status (e.g., ``Buyers: 0''; ``almost no market value''). \\
Failure Inference & Three-item measure of how far the launch itself is judged a failure once the displayed outcome has been seen: whether it appears regrettable, a bad decision, or a serious problem for the company (adapted from Bonifield and Cole, 2007). A judgment about the launch, not an inference about what caused the outcome, and not a measure of classic causal-attribution dimensions. \\
Loss of Brand Luxuriousness & Three-item measure of the perceived decline in the parent brand's luxuriousness immediately after exposure (adapted from Hagtvedt and Patrick, 2008). Not a measure of long-term brand equity. \\
\bottomrule
\end{tabular}
\endgroup
\end{table}

\section{Theoretical Background and Hypotheses}

\subsection{Attribution Theory in Consumer Research}

% T1: Attributions about firm actions
Attribution theory proposes that consumers use observable actions to infer unobservable causes, motives, and intentions. Early consumer research showed that promotional messages and salesperson behavior are processed not only for their surface content but also for why the marketer acted (Smith and Hunt, 1978; Mizerski \emph{et al.}, 1979; Friestad and Wright, 1994; Campbell and Kirmani, 2000). Pricing actions receive the same treatment: consumers' responses to a price change depend on the motive they infer for it, and on the source to which the pricing action is attributed (Campbell, 1999, 2007). Price and distribution information may therefore communicate more than the amount consumers are asked to pay.

% T2: Attributions about unfavorable outcomes
Attribution processes also shape consumer responses after products, services, or firms produce unfavorable outcomes. Research on product failure, service failure, and product-harm crises shows that judgments about cause, controllability, and responsibility change anger, complaints, repurchase, and blame (Folkes, 1984; Folkes \emph{et al.}, 1987; Weiner, 2000; Dawar and Pillutla, 2000; Lei \emph{et al.}, 2012). Thus, an unfavorable outcome affects the firm partly through the meaning consumers assign to that outcome.

% T3: Connecting inferences about actions with evaluations of later outcomes
Research on inferences from marketer actions and research on responses to unfavorable outcomes address different judgments, but both can arise in the same marketplace episode: initial price or distribution information is presented first, and buyer-count and market-value information is displayed later. Related work shows that prior marketing information can shape how consumers interpret subsequent evidence, particularly when that evidence is ambiguous (Deighton, 1984; Hoch and Ha, 1986; Ha and Hoch, 1989). Consumers need not evaluate the later information independently of the earlier information: initial price or distribution information may affect the basis on which consumers interpret subsequently displayed buyer and market-value information. Attribution theory therefore provides the organizing framework for the predicted sequence in this research.

% T3b: The two-stage architecture inherited from crisis research
The product-harm and brand-crisis literature has built a two-stage architecture on this insight: a firm action taken before any adversity --- corporate social responsibility, the expectations a brand has set, or base-rate information --- changes how consumers interpret a subsequent negative event and how far it damages the brand (Dawar and Pillutla, 2000; Klein and Dawar, 2004; Lei \emph{et al.}, 2012), and recent work extends the same architecture to prior brand positioning such as activism, whose protection in one crisis can become a penalty in another (Francioni \emph{et al.}, 2026). This research adopts that architecture for a setting it has not addressed: the prior action is neither a virtue signal nor a set expectation but the launch terms themselves, and the subsequent event is not a defect or scandal but publicly displayed market outcomes. The empirical program therefore proceeds in the order of the episode: Studies 1 and 2 concern judgments formed from the launch terms alone, and Studies 3a--3c concern the interpretation of the outcome information that follows.

% T4: Bridge to brand extension feedback effects
Research on brand extension feedback effects shows that information about an extension can alter evaluations of the parent brand, including dilution following negative extension information (G\"urhan-Canli and Maheswaran, 1998; Ahluwalia and G\"urhan-Canli, 2000; Milberg \emph{et al.}, 2023). The accessibility--diagnosticity perspective explains why extension information that is accessible and perceived as relevant may be used in parent-brand judgments (Ahluwalia and G\"urhan-Canli, 2000). Our focal outcome is Loss of Brand Luxuriousness, an immediate post-exposure perception in the domain of perceived brand luxury (Vigneron and Johnson, 2004; Hagtvedt and Patrick, 2008). We first consider responses to initial price and distribution information before turning to Failure Inference after buyer and market-value information is displayed.

\subsection{Consumer Responses to Initial Price and Distribution Information}

% T5: Positive price and price-quality inference
A positive price presents the extension as a product offered in exchange for payment. A long research tradition shows that price supports quality inferences, and that low prices can shift value perceptions downward, a concern that should be particularly acute for luxury offerings, whose value rests substantially on prestige (Rao and Monroe, 1989; Dodds \emph{et al.}, 1991). Consistent with price--quality inference, a low positive price may lead consumers to evaluate the paid digital extension as less luxurious.

% T6: The free condition in Study 1
A zero price is not merely a smaller positive price; research on the zero-price effect shows that free offers can be evaluated discontinuously from low-priced offers (Shampanier \emph{et al.}, 2007). In particular, a low discounted price serves as a natural reference point for estimating value, whereas a free offer is less likely to be evaluated against that reference price (Palmeira and Srivastava, 2013), and reactions to free offers depend on the attributions consumers form about why the firm makes them (Dallas and Morwitz, 2018). In this research the zero-price effect speaks to Study 1's low-versus-free contrast. The free condition in Study 1 therefore need not produce a more negative valuation than the low-price condition.

% T6b: The exclusivity-based rival prediction and the directional argument
A rarity-based account points in the opposite direction. Perceived brand luxury rests in part on exclusivity and restricted access (Vigneron and Johnson, 2004), and luxury goods signal status partly because few can obtain them (Han \emph{et al.}, 2010); from that perspective, free distribution widens access most and should damage perceived luxuriousness at least as much as any positive price. The price-based account predicts otherwise: a low positive price keeps the extension inside a market exchange and actively certifies a low market worth for an object carrying the brand's name, whereas free distribution removes the offering from that price frame, so no low worth is certified even as access widens. H1 states the price-based prediction, and the exclusivity-based alternative makes Study 1 an informative test rather than a foregone conclusion.

% T7: Application to the digital brand extension, H1
Study 1 applies the distinction between the low-price and free conditions to a digital brand extension before any buyer or market-value information is shown. The low-price condition presents the extension as a paid offer priced substantially below other brands' NFTs, and the free condition presents it as requiring no payment; the focal outcome is Loss of Brand Luxuriousness.

\begin{quote}
\textbf{H1:} Before any buyer-count or market-value information is displayed, the low-price condition will produce higher scores on Loss of Brand Luxuriousness than the free condition.
\end{quote}

% T8: Study 2's concrete conditions and product form
The effect of the complimentary-with-a-fashion-magazine condition, relative to the comparable-price condition, may depend on whether the offering is physical or digital. Physical luxury goods carry cues of costly materials, craftsmanship, and production sacrifice, whereas digital goods make those cues less available and tend to be valued lower than physical equivalents (Atasoy and Morewedge, 2018; Bao \emph{et al.}, 2025; Lu \emph{et al.}, 2025). Product form is examined as a moderator of this concrete comparable-price versus magazine-gift contrast.

% T9: The physical offering
For a physical luxury item, presentation as a fashion-magazine gift may conflict with expectations of costly materials, craftsmanship, and selective exchange: the same object is now positioned as a promotional giveaway inside another product's transaction. This conflict is expected to appear first as a higher score on Loss of Product Luxuriousness.

% T10: The digital offering
For a digital offering, material production cues are less available as a basis for evaluating luxuriousness, and the value of digital collectibles rests more on brand association and social meaning than on inferred production sacrifice (Atasoy and Morewedge, 2018; Hofstetter \emph{et al.}, 2024). The same gift presentation therefore has less to conflict with, and the difference in Loss of Product Luxuriousness between the two Study 2 conditions is expected to be smaller for the digital than for the physical offering.

% T11: From product evaluation to parent-brand evaluation, H2
Because the offering carries the parent brand, a higher score on Loss of Product Luxuriousness may be accompanied by a higher score on Loss of Brand Luxuriousness, consistent with brand extension feedback effects. We therefore estimate a moderated indirect effect with Loss of Product Luxuriousness as the mediator, allowing for a residual direct effect.

\begin{quote}
\textbf{H2:} The complimentary-with-a-fashion-magazine condition, relative to the comparable-price condition, will increase Loss of Product Luxuriousness more for a physical than for a digital offering; accordingly, the estimated indirect effect on Loss of Brand Luxuriousness through Loss of Product Luxuriousness will be stronger for the physical offering.
\end{quote}

\subsection{Displayed Buyer and Market-Value Information and Failure Inference}

% T12: Distinguishing displayed outcomes from Failure Inference
Displayed information showing ``Buyers: 0'' and almost no market value is unfavorable, but it does not by itself establish that the initiative failed. The valence of the displayed information is a property of the stimulus; Failure Inference is an additional judgment about what that information implies about the launch. Consumers must still decide whether those displayed outcomes imply that the launch itself was regrettable or poorly judged.

% T13: Purchased/comparable versus complimentary magazine gift
The earlier purchased/comparable versus complimentary magazine-gift condition provides context for interpreting the same displayed buyer and market-value information. The purchased/comparable condition presents the initial distribution as a paid market exchange, which makes buyer activity a readily available standard for judging the launch. The complimentary magazine-gift condition does not present the initial distribution as the same paid exchange, so the absence of buyers alone may provide a weaker basis for judging the launch as a failure. Failure Inference should therefore be higher after the purchased/comparable condition than after the complimentary magazine-gift condition.

% T14: Definition and limits of Failure Inference
Failure Inference is a three-item measure of how far the launch itself is judged a failure once the displayed outcome has been seen: whether it appears regrettable, a bad decision, or a serious problem for the company (adapted from Bonifield and Cole, 2007; item wording in the Appendix). It records a judgment about the launch rather than an inference about what caused the displayed outcome, and it is not a measure of the classic causal-attribution dimensions of locus, stability, controllability, or responsibility (Folkes, 1984; Weiner, 2000). Because the items evaluate a launch made in the brand's name, higher Failure Inference should be associated with higher scores on Loss of Brand Luxuriousness.

% T15: Link to Loss of Brand Luxuriousness, H3a
If the complimentary magazine-gift condition lowers Failure Inference, and Failure Inference in turn accompanies higher Loss of Brand Luxuriousness immediately after exposure, then the condition should carry a negative estimated indirect effect on Loss of Brand Luxuriousness through Failure Inference.

\begin{quote}
\textbf{H3a:} Under information showing Buyers = 0 and almost no market value, the complimentary magazine-gift condition, relative to the purchased/comparable condition, will produce lower scores on Failure Inference and a negative estimated indirect effect on Loss of Brand Luxuriousness through Failure Inference.
\end{quote}

\subsection{Specific Presentation Conditions in Studies 3b and 3c}

% T16: Two contrasts, reported separately
Studies 3b and 3c test whether the H3a pattern differs across two separate presentation contrasts: an image dominated by a central GG logo versus a bee image with a small logo (Study 3b), and an official brand website versus a news article (Study 3c). Because the two contrasts manipulate different features of the presentation, they are theorized and reported separately.

% T17: Central GG logo versus bee image with a small logo, H3b
Study 3b compares an image dominated by a central GG logo with an image dominated by a bee motif that contains a smaller logo. Brand prominence research, which defines prominence as the conspicuousness of brand marks on a product (Han \emph{et al.}, 2010), provides one interpretive cue: an extension that foregrounds the brand's most recognizable identifier may be harder to evaluate apart from the brand itself, so the displayed outcome of that extension may retain its evaluative force regardless of how the launch was priced or distributed. The prediction is stated for the two images actually used (stimulus texts in Supplementary Information A-4).

\begin{quote}
\textbf{H3b:} The reduction in Failure Inference for the complimentary fan-book-gift condition relative to the comparable-price condition, and the corresponding negative estimated indirect effect on Loss of Brand Luxuriousness through Failure Inference, will be closer to zero for the image with a central GG logo than for the bee image with a small logo.
\end{quote}

The direction of H3b was preregistered (AsPredicted \#259899) under different construct labels: the preregistration predicted that a flagship-logo NFT would amplify the effect of a free offer on loss of perceived brand luxury through anticipated firm regret, the moderation stated here with the moderator carried as the implemented image contrast and the mediating judgment as Failure Inference.

% T18: Official brand website versus news article, H3c
Study 3c presents the launch and the subsequent buyer and market-value information either on an official brand website or in a news article, holding the described offering constant. On the official brand website, the brand itself is the visible author of both the launch and the reported outcome; in the news article, a third party reports the same launch and outcome. When the brand visibly authors both, responsibility for the initiative remains salient whether the offering was initially free or paid, so the transaction frame should become less diagnostic of whether zero buyers imply a failed launch. Source effects on the interpretation of marketer actions (Campbell, 2007) motivate the contrast, and the prediction is stated for the contrast between the official brand website and the news article.

\begin{quote}
\textbf{H3c:} The reduction in Failure Inference for the free condition relative to the comparable-price condition, and the corresponding negative estimated indirect effect on Loss of Brand Luxuriousness through Failure Inference, will be closer to zero on the official brand website than in the news article.
\end{quote}

The direction of H3c was likewise preregistered (AsPredicted \#268386), phrased as flagship framing --- an official announcement rather than a third-party news report --- weakening the buffering effect of a free offer on loss of perceived brand luxury through perceived brand failure; this article names the implemented contrast by its source and the mediating judgment as Failure Inference.

% T19: A pathway not captured by Failure Inference
The contrast between the official brand website and the news article may also be associated with Loss of Brand Luxuriousness through a pathway not captured by Failure Inference. The formal prediction concerns the conditional indirect effect through Failure Inference; direct and total effects are evaluated separately.

\section{Overview of the Empirical Program}

% O1: Cumulative order
The empirical program proceeds cumulatively. The Preliminary Field Study describes associations between a Free/Paid classification of brand NFT collections and observed marketplace indicators on OpenSea, and between that classification and the VADER compound scores of X posts retrieved with collection-specific queries. Study 1 tests H1 with four price conditions and no subsequent buyer or price information. Study 2 tests H2 in a 2 $\times$ 2 design crossing physical/digital product form with the comparable-price versus complimentary-with-a-fashion-magazine contrast. Study 3a tests H3a after identical information showing Buyers = 0 and almost no market value. Study 3b tests H3b with the central-GG-logo versus bee-image contrast, and Study 3c tests H3c with the official-brand-website versus news-article contrast. A further preregistered experiment presenting favorable buyer and price information (Study 4) is reported in the Supplementary Information and summarized in the Cross-Study Synthesis.

% O2: Roles of the evidence (single home of the mediator scope statement)
The two field analyses are observational: they describe associations at the collection and post level and provide no causal test of H1--H3c. The experiments provide randomized tests of condition effects on perceptual outcomes measured immediately after exposure; in all mediation models the mediators were measured rather than randomized, so we describe indirect effects throughout as process-consistent estimates and report total, direct, and indirect effects in full. Three features of the program strengthen the inferences it can support. First, Studies 3a--3c hold the displayed buyer and market-value information literally identical across conditions, so any condition difference in Failure Inference reflects the interpretive context created by the launch information, not the displayed outcome itself. Second, the first-stage H3a pattern is tested across three different free or complimentary implementations, two brands, and image- and text-based materials, providing conceptual replication rather than repetition of a single stimulus; the implementations remain distinct conditions (Table II; rationale in Supplementary Information H). Third, null results are probed with equivalence tests rather than accepted by default.

% O3: Transparency
Studies 2, 3a, 3b, 3c, and 4 were preregistered on AsPredicted (\#259509, \#265490, \#259899, \#268386, and \#262807, respectively); Study 1 was not preregistered. One deviation from a preregistration is disclosed in the Study 3c method. All studies used general-consumer samples with an interest in NFTs recruited through a crowdsourcing platform; participants read a brief explanation of NFTs before the manipulations. Each study embedded one or more instructed-response checks, and participants who failed any of them were excluded; the checks administered in each study and the resulting exclusions are listed in Supplementary Information H-5. All studies were anonymous online surveys of adults who took part voluntarily and consented before participation. In accordance with the research-conducting institution's ethics-review flowchart, the research was confirmed to comply with the institution's regulations for research involving human participants. Stimuli in Studies 3a--3c and the supplementary study present order-book-style buyer and price displays that approximate the marketplace pages on which such information actually appears. Scale items and reliabilities appear in the Appendix (Table A.1); full scenario texts appear in the Supplementary Information. The measures that enter the mediation models were also examined for discriminant validity: in every study a two-factor solution fitted better than a one-factor alternative, and both the Fornell and Larcker (1981) criterion and the heterotrait--monotrait ratio (Henseler \emph{et al.}, 2015) were satisfied (Supplementary Information H-4). Figure 1 summarizes the design of the empirical program.

\begin{table}[H]
\centering
\caption{Initial Price and Distribution Conditions as Implemented in Each Study}
\begingroup
\doublespacing
\renewcommand{\arraystretch}{1.0}
\begin{tabular}{@{}
  >{\raggedright\arraybackslash}p{(\columnwidth - 4\tabcolsep) * \real{0.14}}
  >{\raggedright\arraybackslash}p{(\columnwidth - 4\tabcolsep) * \real{0.43}}
  >{\raggedright\arraybackslash}p{(\columnwidth - 4\tabcolsep) * \real{0.43}}@{}}
\toprule
\textbf{Study} & \textbf{Paid condition(s) as implemented} & \textbf{Free/complimentary condition as implemented} \\
\midrule
Study 1 & Price higher than / comparable to / substantially lower than other brands' NFTs (three conditions) & ``This NFT will be distributed for free.'' \\
Study 2 & Sold at a price comparable to other luxury brands (physical: regular retail; digital: comparable to other luxury-brand NFTs) & ``Offered as a complimentary gift with a fashion magazine'' (physical and digital versions) \\
Study 3a & Same comparable-price wording as Study 2 (digital only) & Same complimentary magazine-gift wording as Study 2 (digital only) \\
Study 3b & Price comparable to other luxury brands' NFTs & ``Free (complimentary with a fan book)'' \\
Study 3c & Price comparable to other luxury brands' NFTs & ``Distributed for free'' \\
Study 4 (SI) & Price comparable to other luxury brands' NFTs & ``Sale price: free'' \\
\bottomrule
\end{tabular}
\endgroup
\end{table}

\begin{figure}[H]
\centering
\includegraphics[width=\linewidth]{figures/figure_1_overview.png}
\caption{Overview of the Preliminary Field Study and the Experiments}
\end{figure}

\section{Preliminary Field Study}

% F1: Purpose and limits
The field analyses describe two associations: between a Free/Paid classification of brand NFT collections and observed marketplace outcomes on OpenSea, and between that classification and the VADER compound score of each X post retrieved with an NFT-collection-specific query. We do not relabel these scores as consumer sentiment, brand evaluation, or demand, and neither analysis provides a causal test of H1--H3c.

% F2: OpenSea method
From a raw export of 6,209 OpenSea collections retrieved in April 2025, 77 brand-related collections were identified, of which 65 collections tied to 32 focal brands were retained. Each collection was classified from its mint price, that is, the price at which a token was first acquired from the brand on the primary market: MintPrice\_ETH = 0, including complimentary distribution and Forge/Burn origins, as Free (\emph{n} = 37), and MintPrice\_ETH \textgreater{} 0 as Paid (\emph{n} = 28; mean mint price 1.44 ETH). The outcome was the collection's floor price, that is, the lowest price at which any token in the collection was listed for sale on the secondary market at the time of retrieval. Floor price was log-transformed and modeled with a mixed-effects regression with brand as a random intercept, with the classification dummy coded 1 for Free and 0 for Paid; nonparametric and within-brand comparisons served as robustness checks.

% F3: OpenSea results
Median floor prices were similar for Free (0.014 ETH) and Paid (0.017 ETH) collections, and the Free coefficient in the mixed-effects model was not significant ($\beta$ = -0.22, \emph{SE} = 0.42, \emph{z} = -0.53, \emph{p} = .598; Mann--Whitney \emph{U} = 530.5, \emph{p} = .874). A within-brand comparison across the six brands with both Free and Paid collections likewise showed no difference (paired \emph{t} = 0.06, \emph{p} = .955). The Free group's higher mean (0.499 ETH vs.\ 0.099 ETH) reflects outliers, and 26 of the 28 Paid collections traded below their mint price. These null results indicate no detectable disadvantage of free distribution in observed floor prices in this cross-section, without establishing equivalence (details in Supplementary Information D).

% F4: X method
The X data concern posts about specific NFT collections, not generic posts about the parent brands. For each collection, a query combined collection and project names and hashtags with NFT/Web3 disambiguation terms (the query construction procedure is described in Supplementary Information F); each retrieved post was mapped to its collection slug and inherited that collection's Free/Paid classification. Retrieval covered February 9, 2022 through April 28, 2025 and returned 33,530 posts, of which 27,467 were English-language. Sequential cleaning removed duplicate post IDs, accounts with ``bot'' in the username (predominantly mechanical sale and listing notifications), duplicate texts, URL-only posts, very short posts, and brands with fewer than ten posts, yielding 22,841 posts across 32 brands (Free 10,381; Paid 12,460). Official brand and project accounts were not separately excluded, a point we return to below. Each post received a VADER compound score (Hutto and Gilbert, 2014; implementation consistent with Costa \emph{et al.}, 2019). The main model was a linear mixed model with brand as a random intercept and calendar time, measured in months since the earliest retrieved post, as a covariate; brands were classified as Luxury (\emph{n} = 10) or Non-Luxury (\emph{n} = 22; classification in Supplementary Information E).

% F5: X results
The Paid $\times$ Luxury interaction was significant ($\beta$ = -0.41, \emph{SE} = 0.12, \emph{z} = -3.48, \emph{p} \textless{} .001): Paid classification was associated with lower VADER compound scores among Luxury brands ($\beta$ = -0.44, \emph{SE} = 0.12, \emph{z} = -3.76, \emph{p} \textless{} .001) and, far more weakly, among Non-Luxury brands ($\beta$ = -0.03, \emph{SE} = 0.01, \emph{z} = -3.11, \emph{p} = .002). A generalized estimating equations model with negative posts (compound $\leq$ -0.05) as the outcome yielded an odds ratio of 1.24 for Paid (95\% CI {[}1.01, 1.51{]}, \emph{p} = .041), a pattern consistent with a higher prevalence of negative posts; we did not estimate a corresponding model for positive posts. Intraclass correlation indicated substantial between-brand variance (ICC = .245).

% F6: Limitations of the field data (detail in SI I)
These descriptive results are bounded by the data's provenance: the posts are public commentary retrieved by collection-specific queries without wallet linkage, official brand and project accounts were not separately excluded, and the models cluster posts by brand although price is classified at the collection level (full pipeline and audit status in Supplementary Information H). The Luxury interaction rests on few luxury brands with Free collections and is exploratory.

% F7: Transition
These observational associations motivate randomized tests of how initial price and distribution information affects evaluations of a digital brand extension, beginning with initial price before any displayed outcome information.

\section{Study 1}

\subsection{Method}

% S1-1 / S1-2
Study 1 tests H1 using a one-factor, four-level between-subjects design manipulating the launch price of a Gucci NFT (high / comparable / low / free), with no subsequent buyer-count or market-value information. Participants were 674 consumers with an interest in NFTs recruited through a crowdsourcing platform; four who failed the attention check were excluded (\emph{N} = 670; 61\% male; ages 20--69; median age bracket 35--44). Each condition described the launch in one sentence, varying only the price information: sold ``at a price higher than those of NFTs from various other brands,'' ``comparable to those of NFTs from other luxury brands,'' ``substantially lower than those of NFTs from various other brands,'' or ``distributed for free'' (full texts in Supplementary Information A-1). The primary outcome was Loss of Brand Luxuriousness; Behavioral Intention was measured as a separate outcome.

\subsection{Results and Discussion}

% S1-3: H1 test first
The planned H1 contrast was supported: the low-price condition (\emph{M} = 4.54, \emph{SE} = 0.11) produced higher Loss of Brand Luxuriousness than the free condition (\emph{M} = 3.91, \emph{SE} = 0.11; Bonferroni-adjusted \emph{p} \textless{} .001).

% S1-4: Full pattern
The omnibus effect of price condition was significant (\emph{F}(3, 666) = 35.17, \emph{p} \textless{} .001, partial $\eta^2$ = .14), with condition means of 2.93 (high), 3.69 (comparable), 4.54 (low), and 3.91 (free; all \emph{SE}s = 0.11). The low-price condition exceeded every other condition (\emph{p}s \textless{} .001), and the free condition did not differ significantly from the comparable-price condition (\emph{p} = .167). This non-significant difference is not evidence of equivalence: a two one-sided tests procedure with a smallest effect size of interest of \emph{d} = $\pm$0.30 did not establish equivalence between the free and comparable conditions (observed \emph{d} = 0.14; full equivalence and sensitivity analyses in Supplementary Information H). An ANCOVA controlling for self--brand connection and purchase experience left the condition effect unchanged (Supplementary Information B).

% S1-5: Behavioral intention as a separate outcome
Behavioral Intention also varied by condition (\emph{F}(3, 666) = 70.30, \emph{p} \textless{} .001, partial $\eta^2$ = .24): the free condition produced the highest intention (\emph{M} = 5.12, \emph{SE} = 0.11) and the high-price condition the lowest (\emph{M} = 2.85, \emph{SE} = 0.11). We report this as a separate outcome and do not combine it with the brand outcome into a single claim (Figure 2).

% S1-6: Discussion
Study 1 supports H1: a low positive price and free distribution are not interchangeable points on a single price continuum for a luxury digital brand extension. The results do not show that free distribution and comparable pricing are equivalent, and they say nothing about protection of the brand over time. Study 1 was also not preregistered, and the later preregistered studies implemented complimentary conditions rather than plain free distribution, so the low-versus-free contrast rests on this single large-sample test. Because product form may change what price and distribution information implies about an offering, Study 2 turns to the contrast between a physical and a digital version of the same offering.

\begin{figure}[H]
\centering
\includegraphics[width=\linewidth]{figures/figure_2_study1.pdf}
\caption{Study 1: Loss of Brand Luxuriousness (A) and Behavioral Intention (B) by Price Condition}
\end{figure}

\section{Study 2}

\subsection{Method}

% S2-1 / S2-2
Study 2 tests H2 in a 2 (product form: physical vs.\ digital) $\times$ 2 (condition: comparable price vs.\ complimentary with a fashion magazine) between-subjects design, preregistered on AsPredicted (\#259509). Participants were 283 consumers with an interest in NFTs; seven who failed one or more of the three instructed-response checks were excluded (\emph{N} = 276; 58.0\% male; median age bracket 40--44). The stimulus was a Gucci mini-bag, presented either as a physical product or as an NFT image with identical designs. The comparable-price condition described a regular sale at a price comparable to other luxury brands; the complimentary condition described the same offering ``as a complimentary gift with a fashion magazine'' (full texts in Supplementary Information A-2). This complimentary condition changes price, gift status, and distribution context together, and we analyze it as that composite condition.

% S2-3: Measures
The mediator was Loss of Product Luxuriousness and the outcome was Loss of Brand Luxuriousness; Inferred Cost was measured as a manipulation check of the product-form difference (scales in the Appendix).

\subsection{Results and Discussion}

% S2-4: Manipulation check
Physical products were inferred to involve higher costs (\emph{M} = 4.69, \emph{SD} = 1.32) than the digital version (\emph{M} = 3.41, \emph{SD} = 1.62; Welch's \emph{t}(242.14) = 7.12, \emph{p} \textless{} .001, \emph{d} = 0.87), consistent with the intended product-form difference.

% S2-5: Interaction
In a moderated mediation model (PROCESS Model 7; Hayes, 2022; 5,000 bootstrap resamples) with condition (0 = comparable, 1 = complimentary magazine gift) as the predictor, the condition $\times$ product-form interaction on Loss of Product Luxuriousness was significant ($\Delta R^2$ = .04, \emph{F}(1, 272) = 11.83, \emph{p} \textless{} .001). For the physical offering, the complimentary condition increased Loss of Product Luxuriousness (\emph{b} = 0.91, \emph{SE} = 0.24, \emph{t}(272) = 3.85, \emph{p} \textless{} .001, 95\% CI {[}0.44, 1.37{]}); for the digital offering, it did not (\emph{b} = -0.29, \emph{SE} = 0.26, \emph{t}(272) = -1.13, \emph{p} = .259, 95\% CI {[}-0.79, 0.21{]}).

% S2-6: Moderated mediation
The index of moderated mediation was significant (Index = 0.83, \emph{SE} = 0.26, 95\% CI {[}0.36, 1.35{]}): the estimated indirect effect on Loss of Brand Luxuriousness through Loss of Product Luxuriousness was significant for the physical offering (Effect = 0.63, \emph{SE} = 0.16, 95\% CI {[}0.32, 0.95{]}) but not for the digital offering (Effect = -0.20, \emph{SE} = 0.20, 95\% CI {[}-0.60, 0.17{]}). A simple moderation model without the mediator showed the same interaction on Loss of Brand Luxuriousness ($\Delta R^2$ = .06, \emph{F}(1, 272) = 19.79, \emph{p} \textless{} .001; physical: \emph{b} = 1.41, \emph{SE} = 0.25, \emph{p} \textless{} .001; digital: \emph{b} = -0.24, \emph{SE} = 0.27, \emph{p} = .384). Path estimates appear in Figure 3 and Table III.

% S2-7: Residual direct effect
Controlling for the mediator, the direct effect of condition on Loss of Brand Luxuriousness remained significant (\emph{b} = 0.39, \emph{SE} = 0.15, \emph{t}(273) = 2.63, \emph{p} = .009, 95\% CI {[}0.10, 0.69{]}), indicating an additional pathway not carried by Loss of Product Luxuriousness.

% S2-8: Discussion
Consistent with H2, the concrete magazine-gift versus comparable-price contrast raised Loss of Product Luxuriousness, and thereby Loss of Brand Luxuriousness, for the physical offering but not for its digital counterpart. For the digital offering, then, the initial-evaluation route was muted; if launch terms matter for digital extensions, their effect should surface once outcome information appears and must be interpreted. Studies 3a--3c address that interpretation stage under the same initial condition contrast.

\begin{figure}[H]
\centering
\includegraphics[width=\linewidth]{figures/figure_3_study2.jpg}
\caption{Study 2: Moderated Mediation Model (Unstandardized Coefficients).}
\raggedright\emph{Note.} * \emph{p} \textless{} .05, ** \emph{p} \textless{} .01, *** \emph{p} \textless{} .001, \textsuperscript{ns} = not significant.
\end{figure}

\begin{table}[H]
\centering
\caption{Study 2: Moderated Mediation Results (Summary)}
\begin{tabular}{@{}
  >{\raggedright\arraybackslash}p{(\columnwidth - 6\tabcolsep) * \real{0.3333}}
  >{\raggedright\arraybackslash}p{(\columnwidth - 6\tabcolsep) * \real{0.2222}}
  >{\raggedright\arraybackslash}p{(\columnwidth - 6\tabcolsep) * \real{0.2222}}
  >{\raggedright\arraybackslash}p{(\columnwidth - 6\tabcolsep) * \real{0.2222}}@{}}
\toprule
\textbf{Quantity} & \textbf{Effect} & \textbf{Boot}\emph{SE} & \textbf{95\% CI} \\
\midrule
Index of moderated mediation & 0.83 & 0.26 & {[}0.36, 1.35{]} \\
Indirect effect, digital offering & -0.20 & 0.20 & {[}-0.60, 0.17{]} \\
Indirect effect, physical offering & 0.63 & 0.16 & {[}0.32, 0.95{]} \\
\bottomrule
\end{tabular}
\end{table}

\section{Study 3a}

\subsection{Method}

% S3a-1 / S3a-2
Study 3a tests H3a in a one-factor, two-level between-subjects design (purchased/comparable vs.\ complimentary magazine gift), preregistered on AsPredicted (\#265490). Participants were 140 consumers with an interest in NFTs; five who failed one or both of the two instructed-response checks were excluded (\emph{N} = 135; 61.5\% male; median age bracket 40--44). The stimulus and the two initial conditions were identical to the digital conditions of Study 2. After the launch phase, all participants saw the same information framed as the situation one month later: an order book displaying ``Buyers: 0 (none)'' and the statement that ``at present, this NFT has almost no market value'' (Supplementary Information A-3).

% S3a-3: Measures
The mediator was Failure Inference and the outcome was Loss of Brand Luxuriousness, both three-item measures (items in the Appendix). A single exploratory item, ``I believe this NFT is a product intended to generate revenue for the company'' (seven-point scale), is also reported.

\subsection{Results and Discussion}

% S3a-4: Condition -> Failure Inference
The complimentary magazine-gift condition produced lower Failure Inference than the purchased/comparable condition (\emph{b} = -1.17, \emph{SE} = 0.23, \emph{t}(133) = -5.08, \emph{p} \textless{} .001, 95\% CI {[}-1.63, -0.71{]}; \emph{R}\textsuperscript{2} = .16), supporting the first stage of H3a: under identical displayed information, the two initial conditions were followed by different Failure Inference.

% S3a-5: Failure Inference -> brand outcome
Failure Inference was positively associated with Loss of Brand Luxuriousness (\emph{b} = 0.59, \emph{SE} = 0.08, \emph{t}(132) = 7.11, \emph{p} \textless{} .001, 95\% CI {[}0.42, 0.75{]}).

% S3a-6: Indirect, direct, and total effects
In the mediation model (PROCESS Model 4; Hayes, 2022; 5,000 bootstrap resamples), the estimated indirect effect through Failure Inference was negative and significant (Effect = -0.69, Boot\emph{SE} = 0.18, 95\% CI {[}-1.05, -0.38{]}), the direct effect was not significant (\emph{b} = 0.15, \emph{SE} = 0.24, \emph{t}(132) = 0.65, \emph{p} = .520), and the total effect was negative (\emph{b} = -0.53, \emph{SE} = 0.26, \emph{t}(133) = -2.07, \emph{p} = .040). These results support H3a.

% S3a-7: Single exploratory item
On the single exploratory item, ``I believe this NFT is a product intended to generate revenue for the company,'' the purchased/comparable condition scored higher (\emph{M} = 4.74, \emph{SD} = 1.32) than the complimentary condition (\emph{M} = 3.86, \emph{SD} = 1.68; Welch's \emph{t}(121.33) = 3.38, \emph{p} = .001, \emph{d} = 0.59). This is a single author-developed item, not a validated construct measure, and an exploratory, non-preregistered serial model including it appears in Supplementary Information C.

% S3a-8: Discussion
Study 3a provides process-consistent evidence for H3a: the same displayed information showing zero buyers and almost no market value was followed by lower Failure Inference, and thereby lower Loss of Brand Luxuriousness, when the launch had been presented as a magazine gift rather than as a paid sale. Because the displayed information was literally identical across conditions, the difference is not attributable to the outcome information itself, and is consistent with the initial conditions having shaped how that information was read. Study 3b asks whether this pattern survives when the extension's image foregrounds the brand's central identifier.

\section{Study 3b}

\subsection{Method}

% S3b-1 / S3b-2
Study 3b tests H3b in a 2 (condition: comparable price vs.\ complimentary fan-book gift) $\times$ 2 (image: central GG logo vs.\ bee image with a small logo) between-subjects design, preregistered on AsPredicted (\#259899). Participants were 282 consumers with an interest in NFTs; one who failed the attention check was excluded (\emph{N} = 281; 59.8\% male; median age bracket 40--44). One image was dominated by a central GG logo; the other was dominated by a bee motif drawn from the brand's art portfolio with only a small logo (Supplementary Information A-4). The free condition here read ``Free (complimentary with a fan book)''; the displayed buyer and market-value information was identical to Study 3a.

% S3b-3: Measures and check
Failure Inference and Loss of Brand Luxuriousness were measured as in Study 3a. Participants also completed an ancillary single-item judgment concerning whether the NFT appeared to be a flagship product (Hubert \emph{et al.}, 2017). Because this item does not measure logo prominence, we do not treat it as a validation of the image manipulation; it is reported in Supplementary Information H.

\subsection{Results and Discussion}

% S3b-5: Interaction
In the moderated mediation model (PROCESS Model 7; 5,000 bootstrap resamples; image coded 0 = bee image, 1 = central GG logo), the condition $\times$ image interaction on Failure Inference was significant ($\Delta R^2$ = .02, \emph{F}(1, 277) = 6.23, \emph{p} = .013). For the bee image, the fan-book-gift condition lowered Failure Inference (\emph{b} = -0.60, \emph{SE} = 0.23, \emph{t}(277) = -2.57, \emph{p} = .011; \emph{M} = 3.72 vs.\ 4.31); for the central-GG-logo image, it did not (\emph{b} = 0.23, \emph{SE} = 0.24, \emph{t}(277) = 0.98, \emph{p} = .328; \emph{M} = 4.97 vs.\ 4.73).

% S3b-6: Conditional indirect effects
The index of moderated mediation was significant (Index = 0.45, \emph{SE} = 0.19, 95\% CI {[}0.10, 0.87{]}). The conditional indirect effect on Loss of Brand Luxuriousness was negative for the bee image (Effect = -0.33, \emph{SE} = 0.14, 95\% CI {[}-0.62, -0.08{]}) and near zero for the central-GG-logo image (Effect = 0.13, \emph{SE} = 0.13, 95\% CI {[}-0.12, 0.39{]}); the direct effect was not significant (\emph{b} = 0.10, \emph{SE} = 0.16, \emph{p} = .549). Total effects showed the same divergence ($\Delta R^2$ = .03, \emph{F}(1, 277) = 9.82, \emph{p} = .002): for the bee image the fan-book-gift condition lowered Loss of Brand Luxuriousness (\emph{b} = -0.60, \emph{SE} = 0.26, \emph{p} = .020), whereas for the central-GG-logo image it \emph{raised} it (\emph{b} = 0.54, \emph{SE} = 0.26, \emph{p} = .038). Path estimates appear in Figure 4 and Table IV.

% S3b-7: Discussion
Consistent with H3b, the reduction in Failure Inference and the negative estimated indirect effect were confined to the bee image with a small logo and were absent---and at the level of total effects even reversed---for the image with a central GG logo. Because the two images differ in motif and composition as well as logo prominence, alternative visual explanations remain possible (Supplementary Information H). Study 3c examines a different presentation contrast, the source in which the launch and its outcome are reported.

\begin{figure}[H]
\centering
\includegraphics[width=\linewidth]{figures/figure_4_study3b.jpg}
\caption{Study 3b: Moderated Mediation Model (Unstandardized Coefficients).}
\raggedright\emph{Note.} * \emph{p} \textless{} .05, ** \emph{p} \textless{} .01, *** \emph{p} \textless{} .001, \textsuperscript{ns} = not significant.
\end{figure}

\begin{table}[H]
\centering
\caption{Study 3b: Moderated Mediation Results (Summary; Image Coded 0 = Bee Image, 1 = Central GG Logo)}
\begin{tabular}{@{}
  >{\raggedright\arraybackslash}p{(\columnwidth - 6\tabcolsep) * \real{0.3846}}
  >{\raggedright\arraybackslash}p{(\columnwidth - 6\tabcolsep) * \real{0.2051}}
  >{\raggedright\arraybackslash}p{(\columnwidth - 6\tabcolsep) * \real{0.2051}}
  >{\raggedright\arraybackslash}p{(\columnwidth - 6\tabcolsep) * \real{0.2051}}@{}}
\toprule
\textbf{Quantity} & \textbf{Effect} & \textbf{Boot}\emph{SE} & \textbf{95\% CI} \\
\midrule
Index of moderated mediation & 0.45 & 0.19 & {[}0.10, 0.87{]} \\
Indirect effect, bee image & -0.33 & 0.14 & {[}-0.62, -0.08{]} \\
Indirect effect, central GG logo & 0.13 & 0.13 & {[}-0.12, 0.39{]} \\
\bottomrule
\end{tabular}
\end{table}

\section{Study 3c}

\subsection{Method}

% S3c-1 / S3c-2
Study 3c tests H3c in a 2 (condition: comparable price vs.\ free) $\times$ 2 (source: official brand website vs.\ news article) between-subjects design, preregistered on AsPredicted (\#268386). To vary the setting, the brand was Louis Vuitton and all material was text-only. From 287 consumers with an interest in NFTs, five who failed the attention check were excluded (\emph{N} = 282; 55.3\% male; median age bracket 40--44). In the official-brand-website source, the launch and the one-month-later information appeared on ``louisvuitton.com / LOUIS VUITTON / OFFICIAL''; in the news-article source, the same launch and information appeared in a ``Fashion Tech News'' editorial article. The described offering, its limitation to 1,000 units, and the subsequent information (initial price; ``Buyers: 0 (none)''; ``there are currently no buyers'') were held constant across sources (Supplementary Information A-5).

% S3c-3: Preregistration deviation
One deviation from the preregistration is disclosed here: the preregistration described the free condition as a fan-book bonus (as implemented in Study 3b), whereas the implemented scenario presented plain free distribution (``Distributed for free'').

% S3c-4: Measures and check
Failure Inference and Loss of Brand Luxuriousness were measured as in Studies 3a--3b. A source manipulation-specific check was not administered.

\subsection{Results and Discussion}

% S3c-5: First-stage interaction
In the moderated mediation model (PROCESS Model 7; 5,000 bootstrap resamples; source coded 0 = news article, 1 = official brand website), the condition $\times$ source interaction on Failure Inference was significant ($\Delta R^2$ = .01, \emph{F}(1, 278) = 4.02, \emph{p} = .046). In the news article, the free condition lowered Failure Inference (\emph{b} = -0.95, \emph{SE} = 0.23, \emph{t}(278) = -4.20, \emph{p} \textless{} .001, 95\% CI {[}-1.40, -0.50{]}); on the official brand website, the effect was not significant (\emph{b} = -0.32, \emph{SE} = 0.22, \emph{t}(278) = -1.46, \emph{p} = .145, 95\% CI {[}-0.75, 0.11{]}).

% S3c-6: Conditional indirect effects
The conditional indirect effect on Loss of Brand Luxuriousness was negative in the news article (Effect = -0.76, \emph{SE} = 0.18, 95\% CI {[}-1.14, -0.41{]}) and not significant on the official brand website (Effect = -0.26, \emph{SE} = 0.18, 95\% CI {[}-0.59, 0.10{]}); the index of moderated mediation was significant (Index = 0.51, \emph{SE} = 0.26, 95\% CI {[}0.02, 1.05{]}), supporting H3c's first component.

% S3c-7: Countervailing direct effect
Controlling for Failure Inference, the direct effect of the free condition on Loss of Brand Luxuriousness was \emph{positive} and significant (\emph{b} = 0.38, \emph{SE} = 0.14, \emph{t}(279) = 2.62, \emph{p} = .009, 95\% CI {[}0.09, 0.66{]}). The free condition was thus associated with lower Loss of Brand Luxuriousness through Failure Inference and, simultaneously, with higher Loss of Brand Luxuriousness through a pathway not captured by Failure Inference.

% S3c-8: Total effects and inconsistent mediation
Because the indirect and direct effects run in opposite directions, Study 3c exhibits inconsistent mediation (Hayes, 2022), and total effects differ by source ($\Delta R^2$ = .02, \emph{F}(1, 278) = 4.51, \emph{p} = .035): the total effect of the free condition was negative in the news article (\emph{b} = -0.55, \emph{SE} = 0.27, \emph{p} = .043) and non-significant, with a positive point estimate, on the official brand website (\emph{b} = 0.25, \emph{SE} = 0.26, \emph{p} = .342). Reporting the indirect effect alone would therefore overstate the protective pattern.

% S3c-9: Discussion
Study 3c offers qualified support for H3c: the reduction in Failure Inference and the negative indirect effect were confined to the news article. At the same time, the countervailing positive direct effect indicates that a simple account in which free distribution uniformly reduces brand damage is insufficient; on the official brand website the total effect was, if anything, directionally positive. The mechanism of the direct effect was not measured and is left for future research.

\section{Cross-Study Synthesis}

% C1: Initial evaluation
Two results concern the initial evaluation of the extension before any displayed outcome. In Study 1, a low positive price produced higher Loss of Brand Luxuriousness than plain free distribution, and the free condition's difference from the comparable-price condition was not significant, without establishing equivalence. In Study 2, the complimentary magazine-gift condition raised Loss of Product Luxuriousness, and thereby Loss of Brand Luxuriousness, for the physical version of the offering but not for its digital counterpart, with a residual direct effect left unexplained.

% C2: Failure Inference under identical displayed information
Study 3a shows that identical displayed information---Buyers = 0 and almost no market value---was followed by different Failure Inference depending on the initial condition: lower after the complimentary magazine-gift condition than after the purchased/comparable condition, with a negative estimated indirect effect on Loss of Brand Luxuriousness. The judgment that a launch failed is therefore not fixed by the displayed outcome alone. This first-stage pattern was not tied to one stimulus: it emerged again for the bee image in Study 3b and in the news article in Study 3c, under different implementations, brands, and material formats.

% C3: The two presentation contrasts, separately
The two boundary results are reported separately because they are different manipulations. In Study 3b, the reduction in Failure Inference and the negative indirect effect held for the bee image with a small logo and were absent, and at the level of total effects reversed, for the image with a central GG logo. In Study 3c, the corresponding pattern held in the news article and was absent on the official brand website.

% C4: Competing pathways and non-uniformity
The first-stage pattern recurred under three different implementations of the complimentary or free offering---a magazine gift in Study 3a, a fan-book gift in the bee-image condition of Study 3b, and plain free distribution in the news-article condition of Study 3c---across two brands and across image- and text-based materials. That the same directional result survives these implementation differences is what the design can offer in place of a pooled estimate: convergence across operationalisations rather than one reproducible quantity. Three features of the results bound how far that convergence extends. First, the implemented conditions differ across studies (plain free; magazine gift; fan-book gift; Table II), so the estimates are not interchangeable. Second, Study 3c's inconsistent mediation indicates a countervailing positive direct pathway operating alongside the negative indirect one. Third, in Study 3b's central-GG-logo condition the fan-book-gift condition raised Loss of Brand Luxuriousness at the level of total effects. Any managerial reading must therefore be conditional on the implementation and the presentation context.

% Qualitative cross-study synthesis. The single-paper meta-analysis previously
% reported here was removed: pooling five condition-level indirect effects would
% have combined heterogeneous implemented conditions and two different
% presentation manipulations under one coding, which contradicts the reading
% adopted throughout, and the estimates are not independent.
Across Studies 3a--3c, the first-stage effect was negative in Study 3a, in the bee-image condition, and in the news-article condition, but was close to zero in the central-logo and official-website conditions. Because these estimates arose from different implemented conditions and different presentation manipulations, we treat this pattern as a qualitative cross-study synthesis rather than as evidence of a common pooled effect.

% Study 4 disclosure
Finally, a preregistered supplementary experiment (Study 4, AsPredicted \#262807; \emph{N} = 273) presented favorable displayed information (1,343 buyers and a sharply rising price) after the launch phase, using the bee-image stimulus. Perceived Brand Luxury declined from baseline after the launch information and rose after the favorable information in both the comparable-price and free conditions, with no significant condition differences; the design does not establish differential recovery, nor equivalence between conditions (full report in Supplementary Information G).

\section{General Discussion}

% D1: Opening synthesis
This research separated two judgments that initial price and distribution information can shape. Before any outcome information, that information affected how a luxury digital brand extension was evaluated: a low positive price was costlier for perceived brand luxuriousness than free distribution (Study 1), and a complimentary magazine-gift presentation harmed the perceived luxuriousness of a physical, but not a digital, version of the same offering (Study 2). After identical displayed information showing zero buyers and almost no market value, the same initial information changed Failure Inference: the launch looked less like a failure when it had not been presented as a paid market exchange (Study 3a), except when the extension carried the brand's central identifier (Study 3b) or when launch and outcome appeared on the official brand website (Study 3c), where a countervailing direct pathway also emerged.

\subsection{Theoretical Contributions}

% Three-point contribution statement. Studies 3a-3c carry the primary contribution.
The research makes three theoretical points, and the strongest of them comes from Studies 3a--3c rather than from the price contrast in Study 1. First, it separates outcome valence from failure judgment: information showing no buyers is unfavourable in every condition, yet whether consumers read it as a failed launch varied with how the extension had been launched. Second, it identifies launch terms as interpretive priors: the price and distribution terms a firm chooses before any outcome is observable alter how diagnostic objectively identical market information later appears. Third, it adds a pre-outcome determinant to research on brand extension feedback effects: how far an extension's market outcome travels to the parent brand depends not only on that outcome but on the terms under which the extension was launched.

% D2: Contextual application of attribution theory
Drawing on two field datasets and six experiments, five of them preregistered, the findings illustrate how attribution theory can organize consumer responses to digital brand extensions across a sequence in which firm-chosen launch information precedes publicly displayed market outcomes of the same initiative. Prior work established that consumers infer motives from marketer actions (Campbell, 1999; Dallas and Morwitz, 2018), that prior information shapes the interpretation of later evidence (Deighton, 1984; Hoch and Ha, 1986), and that responses to unfavorable outcomes depend on the judgments formed about them (Folkes, 1984; Weiner, 2000). Our contribution is to connect these elements in one marketplace episode: the same displayed outcome was followed by different Failure Inference depending on the initial price and distribution information.

% D2b: What the two-stage architecture buys
Reading the program against the crisis literature's architecture clarifies what is new. In that literature the prior action that shields a brand is a stock of goodwill or a set expectation --- corporate social responsibility, or an activist stance whose protection can invert across crises (Klein and Dawar, 2004; Francioni \emph{et al.}, 2026) --- and the second-stage judgment allocates responsibility for an acknowledged harm. Here the prior action is a commercial term of the launch itself, carrying no goodwill of its own, and the second-stage judgment sits upstream of blame: whether the displayed outcome amounts to a failure at all. On this reading, launch terms do not buffer the brand by making consumers more forgiving; they change what the outcome information is evidence of. The parallel extends to the boundary results: as activism's protection can become a penalty when the crisis touches values, the complimentary launch's protection vanished --- and at the level of total effects reversed --- when the extension carried the brand's central identifier or the brand itself authored the outcome information.

% D3: Initial price and distribution information
For the initial-evaluation stage, the results refine how launch information matters. The harm associated with a low positive price (Study 1) is consistent with price--quality inference (Rao and Monroe, 1989), and the product-form moderation in Study 2 is consistent with digital offerings providing fewer material cost cues against which a giveaway looks incongruent (Atasoy and Morewedge, 2018); both interpretations are theory-consistent readings of condition effects, not measured mechanisms. The Study 1 result also weighs against a purely exclusivity-based reading of launch terms for immediate perceptions: free distribution widens access most, yet it was the low positive price that carried the highest Loss of Brand Luxuriousness; whether wide free access erodes luxuriousness over longer horizons lies outside our short-run design.

% D4: Failure Inference after displayed information
For the outcome-interpretation stage, the results indicate that the same displayed marketplace information was not read the same way across conditions. Under identical Buyers = 0 information, Failure Inference differed by initial condition, and Failure Inference was strongly associated with the immediate parent-brand outcome. Failure Inference captures whether the launch appears regrettable, ill-judged, or problematic---a judgment about the launch itself rather than an attribution of its cause. The condition differences are unlikely to reduce to the arithmetic of what recipients stood to lose: had they done so, the reduction in Failure Inference should have appeared wherever the launch was complimentary, yet it was absent for the image dominated by the brand's central identifier and on the official brand website, which points to an interpretive process conditioned by how the launch was presented.

% D5: The two presentation contrasts
The boundary evidence should be read at the granularity at which it was produced. The image result (Study 3b) concerns one pair of images that differ in motif, composition, and logo prominence; brand prominence research (Han \emph{et al.}, 2010) offers one reading, but the contrast is composite. The source result (Study 3c) concerns an official brand website versus a news article; visible brand authorship of both launch and outcome is one difference between these sources, but source credibility and related constructs were not measured. Treating either result as a general law of logo prominence or of brand-authored communication would exceed the designs.

% D6: Brand extension feedback effects
For research on brand extension feedback effects (Ahluwalia and G\"urhan-Canli, 2000; Milberg \emph{et al.}, 2023), the results add a pre-outcome determinant of feedback from market performance: how an extension's displayed performance related to parent-brand evaluation depended on how the extension had been launched. The accessibility--diagnosticity perspective offers a coherent frame---launch information may change how diagnostic the displayed outcome seems for judging the initiative---but diagnosticity was not measured, so this remains an organizing interpretation.

% D7: Countervailing direct and indirect pathways
Study 3c's inconsistent mediation should be read as a boundary signal rather than dismissed as a nuisance. The free condition simultaneously predicted lower Loss of Brand Luxuriousness through lower Failure Inference and higher Loss of Brand Luxuriousness through an unmeasured pathway, and on the official brand website the net association was directionally positive. Together with the total-effect reversal in Study 3b's central-GG-logo condition, this indicates that the launch conditions studied here can help and harm the parent brand at once, through different routes. Because suppression patterns estimated with measured mediators can also arise from omitted variables, the countervailing pathway awaits direct measurement.

\subsection{Managerial Implications}

% D8: Implication 1
First, a low positive launch price is not necessarily the safe middle ground it appears to be. In Study 1, the low-price condition was the most damaging for perceived brand luxuriousness---worse than free distribution---while free distribution also produced the highest acquisition intention. Managers weighing a discounted launch for a digital extension should treat ``cheap'' and ``free'' as qualitatively different signals, not adjacent points on one scale.

% D9: Implication 2
Second, initial price and distribution information should be planned jointly with the market information that platforms will later display. Because marketplaces make buyer counts and prices public, the launch decision shapes the frame in which those later displays are read: after a paid launch, zero buyers looked like failure; after a complimentary launch, less so---unless the extension carried the brand's central identifier or the outcome appeared on the brand's own website. These qualifications also mean that no single ``free'' recipe generalizes: plain free distribution, a magazine gift, and a fan-book gift are different implementations, and our evidence ties each result to its implementation.

\subsection{Limitations and Future Research}

% D10: Attribution evidence
The attribution account is deliberately bounded. Inferred motives were captured by one exploratory single item in Study 3a; responsibility and the classic causal-attribution dimensions were not measured; and the mediators (Loss of Product Luxuriousness; Failure Inference) were measured, not randomized, so all indirect effects are process-consistent estimates. Direct measurement of action attributions, and designs that manipulate the mediator, are natural next steps.

% D11: Design and field limitations
Several design features limit generality. The complimentary conditions bundle price, gift status, and distribution context; the Study 3b images differ in more than logo prominence; Study 3c deviated from its preregistration in the free condition's wording and used a check item written for the image contrast; and all experiments measured immediate perceptions of hypothetical launches in one product category. The field data are observational: posts were retrieved by collection-specific queries without wallet linkage, official and project accounts were not separately excluded, and the models cluster by brand although price varies at the collection level.

% D12: Future research
Future work could measure action attributions directly (which motives consumers infer from a paid, discounted, gift, or free launch), manipulate them, and trace their consequences for the standards by which displayed outcomes are judged; could unbundle price, gift, and distribution context factorially; could test longitudinal parent-brand outcomes beyond immediate perceptions; and could model field data at the collection level with official-account auditing.

% D13: Conclusion
\subsection{Conclusion}

Initial price and distribution information shapes how consumers evaluate a digital brand extension and how they interpret subsequently displayed buyer-count and market-value information, with immediate consequences for perceived brand luxuriousness. For luxury brands entering markets whose outcomes are displayed in public, the launch decision is therefore also a decision about how later market information will be read.

\newpage
\section*{Declaration of generative AI and AI-assisted technologies in the manuscript preparation process}

During the preparation of this work, the authors used generative AI and AI-assisted tools to refine wording and improve the English of the manuscript. After using these tools, the authors reviewed and edited the content as needed and take full responsibility for the content of the published article.

\newpage
\section*{References}
\begin{referenceshang}

Ahluwalia, R. and G\"urhan-Canli, Z. (2000), ``The effects of extensions on the family brand name: an accessibility--diagnosticity perspective'', \emph{Journal of Consumer Research}, Vol. 27 No. 3, pp. 371--381.

Atasoy, O. and Morewedge, C.K. (2018), ``Digital goods are valued less than physical goods'', \emph{Journal of Consumer Research}, Vol. 44 No. 6, pp. 1343--1357.

Bao, W., Hudders, L., Yu, S. and Beuckels, E. (2025), ``Virtual luxury in the metaverse: NFT-enabled value recreation in luxury brands'', \emph{International Journal of Research in Marketing}, Vol. 42 No. 3, pp. 557--572, doi: 10.1016/j.ijresmar.2024.01.002.

Bonifield, C. and Cole, C. (2007), ``Affective responses to service failure: anger, regret, and retaliatory versus conciliatory responses'', \emph{Marketing Letters}, Vol. 18 No. 1--2, pp. 85--99, doi: 10.1007/s11002-006-9006-6.

Brandes, L. and D\"olp, C. (2025), ``Non-fungible tokens (NFTs) as digital brand extensions: evidence on financial performance and parent-brand spillovers'', \emph{International Journal of Research in Marketing}, Vol. 42 No. 3, pp. 491--521, doi: 10.1016/j.ijresmar.2025.01.001.

Campbell, M.C. (1999), ``Perceptions of price unfairness: antecedents and consequences'', \emph{Journal of Marketing Research}, Vol. 36 No. 2, pp. 187--199.

Campbell, M.C. (2007), ```Says who?!' How the source of price information and affect influence perceived price (un)fairness'', \emph{Journal of Marketing Research}, Vol. 44 No. 2, pp. 261--271.

Campbell, M.C. and Kirmani, A. (2000), ``Consumers' use of persuasion knowledge: the effects of accessibility and cognitive capacity on perceptions of an influence agent'', \emph{Journal of Consumer Research}, Vol. 27 No. 1, pp. 69--83.

CoinDesk (2023), ``Porsche NFT sales slow after backlash over high price'', available at: \url{https://www.coindesk.com/web3/2023/01/26/porsche-nft-sales-slow-after-backlash-over-high-price/} (accessed 27 February 2026).

Costa, H., Gil, R. and Rita, P. (2019), ``Unfolding the characteristics of incentivized online reviews'', \emph{Journal of Retailing and Consumer Services}, Vol. 52, p. 101836, doi: 10.1016/j.jretconser.2018.12.006.

Dallas, S.K. and Morwitz, V.G. (2018), ``There ain't no such thing as a free lunch: consumers' reactions to pseudo-free offers'', \emph{Journal of Marketing Research}, Vol. 55 No. 6, pp. 900--915, doi: 10.1177/0022243718817010.

Dawar, N. and Pillutla, M.M. (2000), ``Impact of product-harm crises on brand equity: the moderating role of consumer expectations'', \emph{Journal of Marketing Research}, Vol. 37 No. 2, pp. 215--226.

Deighton, J. (1984), ``The interaction of advertising and evidence'', \emph{Journal of Consumer Research}, Vol. 11 No. 3, pp. 763--770.

Dodds, W.B., Monroe, K.B. and Grewal, D. (1991), ``Effects of price, brand, and store information on buyers' product evaluations'', \emph{Journal of Marketing Research}, Vol. 28 No. 3, pp. 307--319, doi: 10.2307/3172866.

Folkes, V.S. (1984), ``Consumer reactions to product failure: an attributional approach'', \emph{Journal of Consumer Research}, Vol. 10 No. 4, pp. 398--409.

Folkes, V.S. (1988), ``Recent attribution research in consumer behavior: a review and new directions'', \emph{Journal of Consumer Research}, Vol. 14 No. 4, pp. 548--565.

Folkes, V.S., Koletsky, S. and Graham, J.L. (1987), ``A field study of causal inferences and consumer reaction: the view from the airport'', \emph{Journal of Consumer Research}, Vol. 13 No. 4, pp. 534--539.

Fornell, C. and Larcker, D.F. (1981), ``Evaluating structural equation models with unobservable variables and measurement error'', \emph{Journal of Marketing Research}, Vol. 18 No. 1, pp. 39--50, doi: 10.1177/002224378101800104.

Francioni, B., De Cicco, R., Curina, I. and Cioppi, M. (2026), ``When brand activism meets adversity: consumer reactions to performance- and value-related crises of varying severity'', \emph{Journal of Product \& Brand Management}, Vol. 35 No. 1, pp. 1--15, doi: 10.1108/JPBM-10-2024-5530.

Friestad, M. and Wright, P. (1994), ``The persuasion knowledge model: how people cope with persuasion attempts'', \emph{Journal of Consumer Research}, Vol. 21 No. 1, pp. 1--31.

Grewal, D., Monroe, K.B. and Krishnan, R. (1998), ``The effects of price-comparison advertising on buyers' perceptions of acquisition value, transaction value, and behavioral intentions'', \emph{Journal of Marketing}, Vol. 62 No. 2, pp. 46--59, doi: 10.1177/002224299806200204.

G\"urhan-Canli, Z. and Maheswaran, D. (1998), ``The effects of extensions on brand name dilution and enhancement'', \emph{Journal of Marketing Research}, Vol. 35 No. 4, pp. 464--473.

Ha, Y.-W. and Hoch, S.J. (1989), ``Ambiguity, processing strategy, and advertising-evidence interactions'', \emph{Journal of Consumer Research}, Vol. 16 No. 3, pp. 354--360.

Hagtvedt, H. and Patrick, V.M. (2008), ``Art infusion: the influence of visual art on the perception and evaluation of consumer products'', \emph{Journal of Marketing Research}, Vol. 45 No. 3, pp. 379--389, doi: 10.1509/jmkr.45.3.379.

Han, Y.J., Nunes, J.C. and Dr\`eze, X. (2010), ``Signaling status with luxury goods: the role of brand prominence'', \emph{Journal of Marketing}, Vol. 74 No. 4, pp. 15--30.

Hayes, A.F. (2022), \emph{Introduction to Mediation, Moderation, and Conditional Process Analysis: A Regression-Based Approach}, 3rd ed., Guilford Press, New York, NY.

Henseler, J., Ringle, C.M. and Sarstedt, M. (2015), ``A new criterion for assessing discriminant validity in variance-based structural equation modeling'', \emph{Journal of the Academy of Marketing Science}, Vol. 43 No. 1, pp. 115--135, doi: 10.1007/s11747-014-0403-8.

Hoch, S.J. and Ha, Y.-W. (1986), ``Consumer learning: advertising and the ambiguity of product experience'', \emph{Journal of Consumer Research}, Vol. 13 No. 2, pp. 221--233.

Hofstetter, R., Fritze, M.P. and Lamberton, C. (2024), ``Beyond scarcity: a social value-based lens for NFT pricing'', \emph{Journal of Consumer Research}, Vol. 51 No. 1, pp. 140--150.

Hubert, M., Florack, A., Gattringer, R., Eberhardt, T., Enkel, E. and Kenning, P. (2017), ``Flag up!\,--- flagship products as important drivers of perceived brand innovativeness'', \emph{Journal of Business Research}, Vol. 71, pp. 154--163, doi: 10.1016/j.jbusres.2016.09.001.

Hutto, C.J. and Gilbert, E. (2014), ``VADER: a parsimonious rule-based model for sentiment analysis of social media text'', \emph{Proceedings of the International AAAI Conference on Web and Social Media}, Vol. 8 No. 1, pp. 216--225, doi: 10.1609/icwsm.v8i1.14550.

Klein, J. and Dawar, N. (2004), ``Corporate social responsibility and consumers' attributions and brand evaluations in a product-harm crisis'', \emph{International Journal of Research in Marketing}, Vol. 21 No. 3, pp. 203--217, doi: 10.1016/j.ijresmar.2003.12.003.

Lei, J., Dawar, N. and G\"urhan-Canli, Z. (2012), ``Base-rate information in consumer attributions of product-harm crises'', \emph{Journal of Marketing Research}, Vol. 49 No. 3, pp. 336--348.

Louis Vuitton (2025), ``Digital Collectibles (US) FAQ'', available at: \url{https://us.louisvuitton.com/eng-us/faq/products/digital-collectibles-us} (accessed 27 February 2026).

Lu, J., Mislavsky, R. and Yang, H. (2025), ``Uncommonly rare: consumers value rarity in digital goods more in larger collections'', \emph{Journal of the Association for Consumer Research}, Vol. 10 No. 2, pp. 206--219, doi: 10.1086/734834.

Milberg, S.J., Cuneo, A., Silva, M. and Goodstein, R.C. (2023), ``Parent brand susceptibility to negative feedback effects from brand extensions: a meta-analysis of experimental consumer findings'', \emph{Journal of Consumer Psychology}, Vol. 33 No. 1, pp. 21--44.

Min, L., Liu, P.J. and Anderson, C.L. (2025), ``The visual complexity = higher production cost lay belief'', \emph{Journal of Consumer Research}, Vol. 51 No. 6, pp. 1167--1185, doi: 10.1093/jcr/ucae044.

Mizerski, R.W., Golden, L.L. and Kernan, J.B. (1979), ``The attribution process in consumer decision making'', \emph{Journal of Consumer Research}, Vol. 6 No. 2, pp. 123--140.

Palmeira, M.M. and Srivastava, J. (2013), ``Free offer $\neq$ cheap product: a selective accessibility account on the valuation of free offers'', \emph{Journal of Consumer Research}, Vol. 40 No. 4, pp. 644--656, doi: 10.1086/671565.

Prada (2026), ``Prada Timecapsule'', available at: \url{https://www.prada.com/us/en/pradasphere/special-projects/2026/prada-timecapsule.html} (accessed 27 February 2026).

Rao, A.R. and Monroe, K.B. (1989), ``The effect of price, brand name, and store name on buyers' perceptions of product quality: an integrative review'', \emph{Journal of Marketing Research}, Vol. 26 No. 3, pp. 351--357, doi: 10.1177/002224378902600309.

Shampanier, K., Mazar, N. and Ariely, D. (2007), ``Zero as a special price: the true value of free products'', \emph{Marketing Science}, Vol. 26 No. 6, pp. 742--757, doi: 10.1287/mksc.1060.0254.

Smith, R.E. and Hunt, S.D. (1978), ``Attributional processes and effects in promotional situations'', \emph{Journal of Consumer Research}, Vol. 5 No. 3, pp. 149--158.

Sung, E., Kwon, O. and Sohn, K. (2023), ``NFT luxury brand marketing in the metaverse: leveraging blockchain-certified NFTs to drive consumer behavior'', \emph{Psychology \& Marketing}, Vol. 40 No. 11, pp. 2306--2325, doi: 10.1002/mar.21854.

The Star (2025), ``Nike sued over closure of crypto business'', available at: \url{https://www.thestar.com.my/tech/tech-news/2025/04/25/nike-sued-over-closure-of-crypto-business} (accessed 27 February 2026).

Vigneron, F. and Johnson, L.W. (2004), ``Measuring perceptions of brand luxury'', \emph{Journal of Brand Management}, Vol. 11 No. 6, pp. 484--506, doi: 10.1057/palgrave.bm.2540194.

Weiner, B. (2000), ``Attributional thoughts about consumer behavior'', \emph{Journal of Consumer Research}, Vol. 27 No. 3, pp. 382--387.

\end{referenceshang}

\newpage
\appendix
\setcounter{table}{0}
\renewcommand{\thetable}{A.\arabic{table}}
\renewcommand{\theHtable}{A.\arabic{table}}
\section{Measurement Scales}

\begin{singlespace}
\begin{longtable}[]{@{}
  >{\raggedright\arraybackslash}p{(\columnwidth - 8\tabcolsep) * \real{0.1665}}
  >{\raggedright\arraybackslash}p{(\columnwidth - 8\tabcolsep) * \real{0.4204}}
  >{\raggedright\arraybackslash}p{(\columnwidth - 8\tabcolsep) * \real{0.1321}}
  >{\raggedright\arraybackslash}p{(\columnwidth - 8\tabcolsep) * \real{0.0939}}
  >{\raggedright\arraybackslash}p{(\columnwidth - 8\tabcolsep) * \real{0.1872}}@{}}
\caption{Summary of Measurement Scales and Cronbach's $\alpha$ Coefficients}\\
\toprule
Measure & Scale Items & Studies & $\alpha$ & Adapted From \\
\midrule
\endfirsthead
\multicolumn{5}{l}{\emph{Table A.1 (continued)}}\\
\toprule
Measure & Scale Items & Studies & $\alpha$ & Adapted From \\
\midrule
\endhead
\bottomrule
\endlastfoot
Loss of Brand Luxuriousness & (1) I feel that this brand has become less luxurious. (2) \ldots less attractive. (3) \ldots less high class. & S1--S3c & .94--.96 & Hagtvedt and Patrick (2008); reverse-worded \\
Loss of Product Luxuriousness & (1) This product seems inexpensive. (2) \ldots less luxurious. (3) \ldots low-class. & S2 only & .85 & Hagtvedt and Patrick (2008); reverse-worded \\
Failure Inference & (1) The company would regret offering this NFT. (2) \ldots bad decision. (3) \ldots serious problem. & S3a--S3c & .87--.92 & Bonifield and Cole (2007) \\
Inferred Cost & (1) I think this product involves substantial material costs. (2) \ldots requires a lot of people and resources. (3) \ldots would be costly to replicate. & S2 only & .91 & Min \emph{et al.} (2025); items developed by authors \\
Behavioral Intention & (1) At this price, I would acquire this NFT. (2) The probability\ldots is high. & S1 only & .97 & Grewal \emph{et al.} (1998) \\
Perceived Brand Luxury & (1) This is a luxurious brand. (2) \ldots prestigious. (3) \ldots high class. & S4 (SI) only & .88 / .92 / .92 & Hagtvedt and Patrick (2008) \\
Single exploratory item (monetization) & (1) I believe this NFT is a product intended to generate revenue for the company. & S3a only & --- & Author-developed single item \\
Ancillary single-item judgment & (1) I think this NFT is the flagship product of the brand. & S3b & --- & Hubert \emph{et al.} (2017); single item; not treated as a manipulation check \\
\end{longtable}
\end{singlespace}
\end{document}
